\begin{itemize}
\item \textbf{提案は AFEW で積を上回った．}差は，4 コーパスの下調べで $+2.38$〜$+2.70$ 点，6 コーパスで $+2.17$〜$+2.92$ 点で，基準を満たしたのは，4 コーパスの下調べで 3 つの K のうち 3 つ，6 コーパスで 3 つだった．
\item \textbf{ほかのコーパスでは差が小さい．}AFEW 以外の差は，4 コーパスの下調べで $+0.01$〜$+1.03$ 点（9 条件），6 コーパスで $+0.09$〜$+0.97$ 点（12 条件）だった．差が負だった条件は 0 個である．基準を満たしたのは，4 コーパスの下調べの MELD の K = 25 だけで，6 コーパスでは 0 条件だった．
\item CREMA-D（声と顔）の差は $+0.38$〜$+0.56$ 点で，3 つの K とも区間は 0 より上だった（3 条件）．差が 1 点に届かないので，基準は満たさない．
\item \textbf{水準を K 本で測り，相性をソースから持ち込む組み合わせは，片方だけの版の損を避けた．}ソースだけで測った表は，MELD で積を下回った（4 コーパスの下調べで $-2.36$〜$-1.90$ 点，6 コーパスで $-2.47$〜$-1.90$ 点．基準を満たして下回ったのは 6 条件）．K 本だけで測った表は，6 コーパスの MOSEI の K = 5（$-0.97$ 点，区間 $-1.8$〜$-0.1$）で区間が 0 より下だった．同じ条件での提案の差は $+0.01$〜$+1.03$ 点で，負の値は無かった．
\item \textbf{「どの感情でどのモダリティが強いか」の成分の寄与は小さい．}提案からこの成分だけを消した版との差は，4 コーパスの下調べで $-0.02$〜$+0.36$ 点，6 コーパスで $-0.15$〜$+0.36$ 点で，基準を満たした条件は 0 個だった．提案と積の差は，主にモダリティごとの重み（水準）から来ている．
\item \textbf{表を全ラベルで測れたとしたときの余地}（実際には使えない参照）：積との差は，AFEW で $+3.55$〜$+4.75$ 点，ほかのコーパスで $+0.62$〜$+1.59$ 点だった．AFEW 以外では，表を全ラベルで測っても，積との差は 1.6 点以下だった．
\item \textbf{提案の答えを疑似ラベルにして融合の分類器を学び直すと，差は AFEW で残った}（6 コーパス）．積を疑似ラベルにした分類器との差は，AFEW で $+1.16$〜$+1.81$ 点（3 つの K とも基準を満たす），ほかのコーパスで $+0.19$〜$+1.14$ 点だった（基準を満たしたのは MELD の K = 25）．
\end{itemize}
